% --------------------------------- Preamble -----------------------------------
\documentclass[11pt, letterpaper, twocolumn]{article}

\usepackage[T1]{fontenc} % use 8-bit fonts
\usepackage[margin=0.75in, columnsep=0.2in, % 3.4 in columns
    headsep=0.28in, footskip=0.42in]{geometry} % paper sizes
\usepackage{amsmath, amsfonts, amssymb} % improved math
\usepackage{bm} % bold math (must be after amsmath)
\usepackage{upquote} % upright single quotes in verbatim
\usepackage{pgfplots} % includes tikz, xcolor, graphicx
\usepackage[colorlinks=true, allcolors=gray]{hyperref}

\raggedbottom % do not force stretching text to fill the page
\setlength{\emergencystretch}{\hsize} % avoid overfull lines
\overfullrule=1mm % black band to show overfull lines
\pdfsuppresswarningpagegroup=1 % remove pdf image warnings
\frenchspacing % consistent spacing between words
\newcommand{\ul}[1] % vectors
    {{}\mkern1mu\underline{\mkern-1mu#1\mkern-1mu}\mkern1mu}
\DeclareSymbolFont{euler}{U}{eur}{m}{n} % symbol font
\DeclareMathSymbol{\PI}{\mathord}{euler}{25} % up pi
\pgfplotsset{compat=newest} % use newest pgfplot version
\newcommand{\EE}[2]{\ensuremath{{{#1}\!\times\!10^{#2}}}}
% ------------------------------------------------------------------------------

\title{Linear Kalman Filter for 2D Motion Tracking}
\author{}
\date{}

\begin{document}
\maketitle

\section*{Description}

This project implements a linear Kalman Filter for estimating two-dimensional
position and velocity. The system dynamics are linear, and the measurement model
consists of direct position measurements observed in a rotated sensor frame.
The project includes both simulation of truth and measurements and recursive 
state estimation using a discrete-time Kalman Filter.

\subsection*{System Model}

The system equations are:
\begin{equation*}
    \dot{\ul{x}} = \textbf{F}\ul{x} + \textbf{B}\ul{u} + \ul{\nu}, \quad \ul{\nu} \sim \mathcal{N}(0,\textbf{Q})
\end{equation*}
\begin{equation*}
    \ul{z} = \textbf{H}\ul{x} + \ul{\eta}, \quad \ul{\eta} \sim \mathcal{N}(0,\textbf{R})
\end{equation*}
\\
where
\begin{align*}
    F &= 
    \begin{bmatrix} 
        0 & 0 & 1 & 0 \\ 
        0 & 0 & 0 & 1 \\
        0 & 0 & 0 & 0 \\
        0 & 0 & 0 & 0 
    \end{bmatrix}
    \quad
    B = 
    \begin{bmatrix} 
        0 & 0 \\ 
        0 & 0 \\
        1 & 0 \\ 
        0 & 1 
    \end{bmatrix}
    \\
    \\
    Q &= 
    \begin{bmatrix} 
        0.01 & 0 & 0 & 0 \\ 
        0 & 0.01 & 0 & 0 \\
        0 & 0 & 0.01 & 0 \\
        0 & 0 & 0 & 0.01 
    \end{bmatrix}
    \quad
    \ul{x} = 
    \begin{bmatrix} 
        p_x \\ 
        p_y \\ 
        v_x \\ 
        v_y 
    \end{bmatrix}
    \\
    \\
    H &= 
    \begin{bmatrix} 
        \cos\theta & \sin\theta & 0 & 0 \\
        -\sin\theta & \cos\theta & 0 & 0 
    \end{bmatrix}
    R = 
    \begin{bmatrix} 
        1 & 0 \\ 
        0 & 0.1 
    \end{bmatrix}
    \ul{u} = 
    \begin{bmatrix} 
        a_x \\ 
        a_y 
    \end{bmatrix}
\end{align*}
and where $(p_x, p_y)$ is the position, $(v_x, v_y)$ is the velocity,
$(a_x, a_y)$ is the input acceleration, and $\theta$ is some arbitrarily
chosen angle. There is a sensor at the origin of the coordinate frame
directly measuring the position states. However, it is rotated by the 
angle $\theta$. The \textbf{H} matrix represents this rotation from the
state's coordinate frame to the sensor's coordinate frame: \\
\includegraphics[width=\linewidth]{../figures/fig_frame.png}

\subsection*{Discretization}

The continuous-time system is converted to a discrete-time model with sampling
period $T = 0.1$ s. The discrete dynamics are written as 
\[\begin{aligned}
\ul{x}_{k+1} &= \Phi\ul{x}_k + \mathbf{B}_d\ul{u}_k + \ul{\nu}_k, &\quad \ul{\nu} &\sim \mathcal{N}(0,\mathbf{Q}_d) \\
\ul{z}_k &= \mathbf{H}\ul{x}_k + \mathbf{D}\ul{u}_k + \ul{\eta}_k, &\quad \ul{\eta} &\sim \mathcal{N}(0,\mathbf{R}) 
\end{aligned}
\]

The matrices $\Phi$, $\mathbf{B}_d$, and $\mathbf{Q}_d$ are computed using 
the van Loan method. This avoids relying on the first-order approximations
and provides a more accurate discretization of both the dynamics and 
process noise.

\section*{Kalman Filter}

State estimation is performed using a standard discrete-time Kalman Filter.

\textbf{Update:}
\[\begin{aligned}
\mathbf{S} &= \mathbf{H}\mathbf{P}_k\mathbf{H}^{\top} + \mathbf{R} \\
\mathbf{K} &= \hat{\mathbf{P}}\mathbf{H}^{\top}\mathbf{S}^{-1} \\
\ul{r} &= \ul{z}_k - \mathbf{H}\hat{\ul{x}}_k - \mathbf{D}\ul{u}_k \\
\hat{\ul{x}}_k &= \hat{\ul{x}}_k + \mathbf{K}\ul{r} \\
\hat{\mathbf{P}}_k &= \hat{\mathbf{P}}_k - \mathbf{K}\mathbf{H}\hat{\mathbf{P}}_k
\end{aligned}
\]

\textbf{Propagate:}
\[\begin{aligned}
\hat{\ul{x}}_{k+1} &= \Phi\hat{\ul{x}}_k + \mathbf{B}_d\ul{u}_k \\
\hat{\mathbf{P}}_{k+1} &= \Phi\hat{\mathbf{P}}_k\Phi^{\top} + \mathbf{Q}_d
\end{aligned}
\]

The initial state estimate is taken from the first truth sample, and the 
initial covariance is chosen as the identity matrix.

\section*{Simulation setup}

Truth data and measurements are generated over a 100 s simulation with sampling
period 0.1 s. The acceleration inputs are defined as
\[
a_x = -sin\left(\frac{2\pi}{5}t\right) \quad a_y = -cos\left(\frac{2\pi}{10}t\right)
\]

The initial state is
\[
\begin{bmatrix}
    0 & 0 & 0.8 & 0
\end{bmatrix}^{\top}
\]

Three cases are studied in the analysis:
\begin{enumerate}
    \item Rotated measurement frame with $\theta = \pi/4$
    \item Unrotated measurement frame with $\theta = 0$
    \item Rotated measurement frame with isotropic measurement covariance $\mathbf{R} = \mathrm{diag}[1.0, 1.0]$
\end{enumerate}

\section*{Results}

Figure~\ref{fig:path} shows the true and estimated two-dimensional paths. 
The Kalman Filter tracks the state closely over the full simulation. \\

\begin{figure}[h]
    \centering
    \includegraphics[width=\linewidth]{../figures/fig_paths.pdf}
    \caption{True and estimated 2D trajectories}
    \label{fig:path}
\end{figure}

\newpage

Figures~\ref{fig:pxerr}--\ref{fig:vyerr} shows the state estimation errors
together with the expected $3-\sigma$ bounds from the filter covariance
and the empirical $3-\sigma$ bounds computed from the full error history. \\

\begin{figure}[h]
    \centering
    \includegraphics[width=\linewidth]{../figures/fig_pxerr.pdf}
    \caption{x-position error}
    \label{fig:pxerr}
\end{figure}

\begin{figure}[h]
    \centering
    \includegraphics[width=\linewidth]{../figures/fig_pyerr.pdf}
    \caption{y-position error}
    \label{fig:pyerr}
\end{figure}

\begin{figure}[h]
    \centering
    \includegraphics[width=\linewidth]{../figures/fig_vxerr.pdf}
    \caption{x-velocity error}
    \label{fig:vxerr}
\end{figure}

\begin{figure}[h]
    \centering
    \includegraphics[width=\linewidth]{../figures/fig_vyerr.pdf}
    \caption{y-velocity error}
    \label{fig:vyerr}
\end{figure}

\newpage

\subsection*{Analysis}

\begin{enumerate}

\item The final state covariance matrix of the estimation with
rotated measurement frame with $\theta = \pi/4$

    \begin{center}
        \begin{tabular}{cccc}
            \hline
            0.049 & 0.033 & 0.02 & 0.011 \\
            0.033 & 0.049 & 0.011 & 0.02 \\
            0.02 & 0.011 & 0.022 & 0.005 \\
            0.011 & 0.02 & 0.005 & 0.022 \\
            \hline
        \end{tabular}
    \end{center}

\item The final state covariance matrix after setting $\theta$ to 0. 

    \begin{center}
        \begin{tabular}{cccc}
            \hline
            0.082 & 0.0 & 0.03 & 0.0 \\
            0.0 & 0.016 & 0.0 & 0.009 \\
            0.03 & 0.0 & 0.027 & 0.0 \\
            0.0 & 0.009 & 0.0 & 0.017 \\
            \hline
        \end{tabular}
    \end{center}
    There are 8 zeros in the new covariance matrix, as opposed to the
    matrix representing a $\pi/4$ radian rotation of the sensor. The
    presence of zeros in the aligned matrix is due to the fact that 
    x-dimension and y-dimension become completely independent problems.
    When there is a rotation of the sensor, we must combine $p_x$ and
    $p_y$ in the H matrix to rotate our estimate into the sensor frame.
    This means that each measurement informs the filter about both
    $p_x$ and $p_y$. When there is no rotation, we essentially get 
    one measurement for $p_x$ and one for $p_y$. When there is a rotation,
    the errors in the x and y states are \textit{coupled}, meaning the
    errors in the x and y states are correlated. This is why there are
    no zeros in our original matrix, and zeros appear where x and y
    dimensions cross in the aligned matrix.

\item The final state covariance matrix after setting $\theta$ back 
    to $\pi/4$ and setting $\bm{R}$ to \texttt{diag[1.0, 1.0]}. 

    \begin{center}
        \begin{tabular}{cccc}
            \hline
            0.082 & 0.0 & 0.03 & 0.0 \\
            0.0 & 0.082 & 0.0 & 0.03 \\
            0.03 & 0.0 & 0.027 & 0.0 \\
            0.0 & 0.03 & 0.0 & 0.027 \\
            \hline
        \end{tabular}
    \end{center}
    There are once again 8 zeros in the covariance matrix, but this time
    appearing for a less intuitive reason (which is actually the dominant
    reason). The zeros reappeared because we restored perfect symmetry to 
    the problem. The rotation of the sensor provides the mechanism to
    couple the x and y dimensions, but it is the \textit{asymmetry} in the
    system that actually creates the correlation. When the measurement noise 
    is asymmetric, the filter trusts one sensor direction more than another,
    forcing the x and y estimation errors to become correlated. By making the
    noise symmetric, any preference in direction is removed, so the filter
    treats both dimensions identically, causing the cross-covariance terms
    between them to resolve to zero.

\end{enumerate}

\section*{Discussion}

The results show that the Kalman Filter provides accurate state estimates while
maintaining a covariance matrix that is consistent with the observed estimation
errors. The comparison between the predicted and empirical $3-\sigma$ bounds 
provides a useful check on whether the filter uncertainty is well calibrated.

The covariance analysis also illustrates how the measurement model affects state
coupling. When the measurement frame is rotated, the position states are mixed
in the sensor model, which produces nonzero off-diagonal covariance terms. When
the measurement frame is aligned with the state coordinates, additional structure
appears in the covariance matrix. Changing the measurement covariance to an 
isotropic form further alters this structure by making the measurement uncertainty
symmetric in the rotated frame.
\end{document}